\chapter[APRESENTAÇÃO DA TRAJETÓRIA DO ALUNO]{APRESENTAÇÃO DA TRAJETÓRIA DO ALUNO}

Atualmente, estou no terceiro semestre do meu curso de engenharia de
software. Antes de começar o curso assisti a um curso de Python do canal Curso em
Vídeo para aprender as sintaxes básicas da linguagem e um início em lógica de
programação. Este curso foi importante pois serviu para aquietar meus principais
medos, achar programação muito chato ou muito difícil. Agora com mais certeza do
que faria entrei na faculdade de Engenharia de Software no primeiro semestre de 2023
e logo no primeiro semestre já participei de uma Hackatona junto com meus amigos
mesmo com o limitado conhecimento que tinha na época, atuando na parte do front-
end, configurando minha primeira experiência com HTML, CSS e JavaScript e
despertando um pouco da minha curiosidade para fazer sites, visto que até o
momento, havia apenas feito programas em back-end.

No final do meu segundo semestre fui aprovado para participar do PET
(Programa de Ensino e Tutoria) Informática da PUCRS. Pelo PET entrei no laboratório
DAVINT e comecei a trabalhar com coleta de dados em Python por meio de uma API
da Google feita para o YouTube.

Quanto à parte do front-end continuei com a curiosidade e sem fazer tantos
projetos, portanto, ao ver a oportunidade de entrar na AGES nesta área abracei a
oportunidade no meu terceiro semestre, passando por diversas dificuldades, mas
conseguindo sempre dar meu máximo e entregar as tarefas previstas.



% Utilize referências também, modificando o arquivo \textit{bibliografia/bibliografia.bib} e citando-os com o comando\cite{artigo}.
    